% TODO Written last, 150 to 250 words, in this order: context, gap, research question and hypothesis, data and network model, method with the null model, main result with a number, what it implies.
% Draft of 2026-10-26: the parts that do not exist yet are marked with \todo.
A model that imitates a musical style has to reproduce rules that bind several voices at once, such as the ban on consecutive fifths. Statistical models of four-part harmony count either pairs of notes, which miss these rules, or whole sonorities, which recur too rarely to cover new music. The configurations in between are what network science calls higher-order interactions, and in the sources we checked no model counts them in a corpus to generate music. We ask whether a random walk on a hypergraph can imitate a musical style, and which parts of the hypergraph it needs. \todo{the hypothesis, once H0 and H1 are fixed} For 352 chorales by Bach and 1,179 movements of Palestrina's masses, set on a common time grid and key, the hyperedges are configurations of a few voices over a few cells, each weighted by its support in the corpus. The walk writes the lower voices under a given melody from these counts alone. Removing one shape of hyperedge at a time, and testing on tunes and masses that the counts have not seen, shows which shapes the imitation needs. \todo{the null model; the main result, with its number; what it implies}
